% 数学ノート集・共通テンプレート v2（積分論2026年度第1部の版面を参考）
% A5判。章内で定義・定理・例・注意等の番号を共通化。
% コンパイル： platex -interaction=nonstopmode math_note_template_v2.tex (2回)
%              dvipdfmx math_note_template_v2.dvi
\documentclass[a5paper,10pt,onecolumn]{jsarticle}
\usepackage[dvipdfmx,top=22mm,bottom=20mm,inner=19mm,outer=18mm,headheight=13pt,headsep=6mm]{geometry}
\usepackage{amsmath,amssymb,amsthm,mathtools}
\usepackage{enumitem,fancyhdr,titlesec}
\usepackage[dvipdfmx,bookmarks=true,colorlinks=false]{hyperref}
\usepackage{pxjahyper}

% ----- 差し替える文書情報 -----
\newcommand{\NoteTitle}{Fourier級数の理論}
\newcommand{\NoteSeries}{数学ノート集}

% ----- 文字・組版 -----
% 参考PDFの質素な明朝＋Computer Modernの印象を引き継ぎながら，
% A5判で窮屈にならないよう本文の行間を少し確保する．
\linespread{1.13}
\setlength{\parindent}{1zw}
\setlength{\parskip}{0pt}
\setlength{\abovedisplayskip}{6pt plus 2pt minus 2pt}
\setlength{\belowdisplayskip}{6pt plus 2pt minus 2pt}
\setlength{\abovedisplayshortskip}{3pt plus 2pt minus 1pt}
\setlength{\belowdisplayshortskip}{4pt plus 2pt minus 1pt}
\setlist[enumerate]{topsep=3pt,itemsep=1pt,parsep=0pt,leftmargin=2.4zw}
\titleformat{\section}{\normalfont\normalsize\bfseries}{\thesection}{0.8zw}{}
\titlespacing*{\section}{0pt}{2.0ex plus .4ex}{.9ex}
\titleformat{\subsection}{\normalfont\normalsize\bfseries}{\thesubsection}{0.7zw}{}
\titlespacing*{\subsection}{0pt}{1.5ex plus .3ex}{.6ex}
\raggedbottom

% ----- 柱とページ番号：ページ上部、罫線なし -----
\fancypagestyle{mathnote}{%
  \fancyhf{}%
  \fancyhead[L]{\footnotesize \NoteSeries\ （\NoteTitle）}%
  \fancyhead[R]{\footnotesize \thepage}%
  \renewcommand{\headrulewidth}{0pt}%
  \renewcommand{\footrulewidth}{0pt}%
}
\pagestyle{mathnote}

% ----- 節内通し番号の定義・定理群 -----
\newtheoremstyle{mathnotetheorem}
 {4.5pt}{4.5pt}{\normalfont}{}%
 {\bfseries}{}{0.65zw}%
 {\thmname{#1}\thmnumber{ #2}\thmnote{（#3）}}
\theoremstyle{mathnotetheorem}
\newtheorem{theorem}{定理}[section]
\newtheorem{definition}[theorem]{定義}
\newtheorem{lemma}[theorem]{補題}
\newtheorem{proposition}[theorem]{命題}
\newtheorem{corollary}[theorem]{系}
\newtheorem{example}[theorem]{例}
\newtheorem{remark}[theorem]{注意}
\renewcommand{\proofname}{証明}
\renewcommand{\qedsymbol}{$\square$}
\makeatletter
\renewenvironment{proof}[1][\proofname]{%
  \par\pushQED{\qed}\normalfont\topsep4\p@\@plus2\p@\relax
  \trivlist\item[\hskip\labelsep\bfseries #1]\ignorespaces
}{\popQED\endtrivlist\@endpefalse}
\makeatother

% ----- 使うノートで適宜拡張する記号 -----
\newcommand{\R}{\mathbb{R}}
\newcommand{\norm}[1]{\lVert #1\rVert}
\newcommand{\abs}[1]{\lvert #1\rvert}

\begin{document}
\section*{読了後にできること}
本稿を読むための前提は，三角関数の直交性，Fourier係数と収束定理．
本稿の具体的な到達課題は，一つの関数のFourier係数を計算して級数を構成する．

\section{周期関数とFourier係数}

本稿では $2\pi$-周期関数を考え，
$[-\pi,\pi]$ 上の積分を用いる．
$L^p(\mathbb T)$ は $2\pi$-周期関数の
Lebesgue可積分同値類の空間を意味する．

\begin{definition}
$f\in L^1(\mathbb T)$ に対して，
整数 $n\in\mathbb Z$ のFourier係数を
\[
 \widehat f(n)=\frac1{2\pi}
 \int_{-\pi}^\pi f(t)e^{-int}\,dt
\]
と定める．
第 $N$ 部分和を
\[
 S_Nf(x)=\sum_{n=-N}^N\widehat f(n)e^{inx}
\]
と定義する．
\end{definition}

\begin{lemma}[複素指数関数系の直交性]\label{lem:orthogonal}
整数 $m,n$ について
\[
 \frac1{2\pi}\int_{-\pi}^\pi e^{imt}\overline{e^{int}}\,dt
 =\delta_{mn}
\]
が成り立つ．
\end{lemma}
\begin{proof}
$m=n$ なら被積分関数は1だから積分は1である．
$m\ne n$ のとき
$\int_{-\pi}^\pi e^{i(m-n)t}dt
 =(e^{i(m-n)\pi}-e^{-i(m-n)\pi})/(i(m-n))=0$
である．
\end{proof}

\begin{proposition}[Besselの不等式]\label{prop:bessel}
$f\in L^2(\mathbb T)$ に対し
\[
 \sum_{n\in\mathbb Z}|\widehat f(n)|^2
 \le\frac1{2\pi}\int_{-\pi}^\pi|f(t)|^2\,dt
\]
が成り立つ．
\end{proposition}
\begin{proof}
$e_n(t)=e^{int}$ と置くと
$\{e_n\}_{n\in\mathbb Z}$ は
内積 $\langle f,g\rangle
 =(2\pi)^{-1}\int f\bar g$ に関して正規直交系である．
有限部分系に対する直交射影のPythagorasの恒等式から
\[
 \sum_{n=-N}^N|\widehat f(n)|^2
 =\norm{S_Nf}_2^2\le\norm f_2^2
\]
である．$N\to\infty$ として結論を得る．
\end{proof}

\section{Riemann--Lebesgueの補題}

\begin{theorem}[Riemann--Lebesgueの補題]\label{thm:rl}
$f\in L^1(\mathbb T)$ なら
\[
 \widehat f(n)\longrightarrow0
 \qquad(|n|\to\infty).
\]
\end{theorem}
\begin{proof}
まず $f=1_{[a,b]}$（$-\pi\le a<b\le\pi$）とすると
$n\ne0$ について
\[
 \widehat f(n)
 =\frac1{2\pi}\frac{e^{-ina}-e^{-inb}}{in}
\]
なので $|\widehat f(n)|\le1/(\pi|n|)\to0$ である．
有限個の区間の指示関数の線形結合についても
線形性から成立する．
一般の $f\in L^1(\mathbb T)$ は，
そのような階段関数 $g$ で
$\norm{f-g}_{L^1}<\varepsilon$ と近似できる．
実際，Lebesgue可測集合を有限個の区間の和で
測度について近似し，単純関数近似と組み合わせればよい．
任意の整数 $n$ に対して
\[
 |\widehat f(n)-\widehat g(n)|
 \le\frac1{2\pi}\int_{-\pi}^\pi|f-g|
 <\frac{\varepsilon}{2\pi}.
\]
$|n|\to\infty$ では $\widehat g(n)\to0$ なので
$\limsup|\widehat f(n)|\le\varepsilon/(2\pi)$ となる．
$\varepsilon>0$ は任意だから結論を得る．
\end{proof}

\section{Dirichlet核と部分和}

\begin{definition}
Dirichlet核を
\[
 D_N(t)=\sum_{n=-N}^Ne^{int}
\]
と定める．
\end{definition}

\begin{lemma}\label{lem:dirichlet-kernel}
$N\ge0$ とすると
\[
 D_N(t)=\frac{\sin((N+\tfrac12)t)}{\sin(t/2)}
 \qquad(t\notin2\pi\mathbb Z)
\]
であり，
\[
 S_Nf(x)=\frac1{2\pi}\int_{-\pi}^\pi f(x-t)D_N(t)\,dt
\]
が成り立つ．
\end{lemma}
\begin{proof}
$e^{it}\ne1$ のとき有限等比級数の和を計算すると
\[
 \sum_{n=-N}^Ne^{int}
 =e^{-iNt}\frac{1-e^{i(2N+1)t}}{1-e^{it}}
 =\frac{\sin((N+\frac12)t)}{\sin(t/2)}
\]
となる．
次に
\begin{align*}
 S_Nf(x)
 &=\sum_{n=-N}^N
   \left(\frac1{2\pi}\int_{-\pi}^\pi
   f(s)e^{-ins}\,ds\right)e^{inx}\\
 &=\frac1{2\pi}\int_{-\pi}^\pi
   f(s)\sum_{n=-N}^Ne^{in(x-s)}\,ds.
\end{align*}
変数変換 $t=x-s$ と周期性から表示式を得る．
\end{proof}

\section{Fej\'er核}

\begin{definition}
Fej\'er核とCes\`aro平均をそれぞれ
\[
 F_N(t)=\frac1{N+1}\sum_{k=0}^ND_k(t),\qquad
 \sigma_Nf=\frac1{N+1}\sum_{k=0}^NS_kf
\]
と定める．
\end{definition}

\begin{lemma}[Fej\'er核の基本性質]\label{lem:fejer}
任意の $N\ge0$ と $t\notin2\pi\mathbb Z$ について
\[
 F_N(t)=\frac1{N+1}
 \left(\frac{\sin((N+1)t/2)}{\sin(t/2)}\right)^2.
\]
さらに
\[
 F_N(t)\ge0,\qquad
 \frac1{2\pi}\int_{-\pi}^\pi F_N(t)\,dt=1,
\]
かつ任意の $\delta\in(0,\pi)$ について
\[
 \int_{\delta\le|t|\le\pi}F_N(t)\,dt\longrightarrow0.
\]
\end{lemma}
\begin{proof}
有限幾何級数を用いると
\[
 \frac1{N+1}\left|\sum_{k=0}^Ne^{ikt}\right|^2
 =\sum_{j=-N}^N
   \left(1-\frac{|j|}{N+1}\right)e^{ijt}.
\]
右辺は
$\frac1{N+1}\sum_{k=0}^ND_k(t)=F_N(t)$ に等しい．
左辺の絶対値平方から非負性と三角関数表示を得る．
各 $D_k$ の定数Fourier係数は1なので，
$\int F_N/(2\pi)=1$ である．
$|t|\ge\delta$ の範囲では
$|\sin(t/2)|\ge\sin(\delta/2)>0$ だから
\[
 0\le F_N(t)\le
 \frac1{(N+1)\sin^2(\delta/2)}.
\]
長さ $2\pi$ 以下の範囲で積分すると
$O((N+1)^{-1})\to0$ を得る．
\end{proof}

\section{Fej\'erの定理と近似}

\begin{theorem}[Fej\'erの定理]\label{thm:fejer}
$f$ が連続な $2\pi$-周期関数なら，
\[
 \norm{\sigma_Nf-f}_\infty\longrightarrow0
\]
が成り立つ．
\end{theorem}
\begin{proof}
核の積分表示から
\[
 \sigma_Nf(x)-f(x)
 =\frac1{2\pi}\int_{-\pi}^\pi
   (f(x-t)-f(x))F_N(t)\,dt
\]
である．
$f$ はコンパクトな円周上で一様連続だから，
任意の $\varepsilon>0$ に対し
$\delta\in(0,\pi)$ が存在して
$|t|<\delta$ なら
$|f(x-t)-f(x)|<\varepsilon$ が全ての $x$ で成り立つ．
従って
\begin{align*}
 |\sigma_Nf(x)-f(x)|
 &\le\frac{\varepsilon}{2\pi}
      \int_{|t|<\delta}F_N(t)\,dt\\
 &\quad+\frac{2\norm f_\infty}{2\pi}
      \int_{\delta\le|t|\le\pi}F_N(t)\,dt\\
 &\le\varepsilon+
 \frac{\norm f_\infty}{\pi}
      \int_{\delta\le|t|\le\pi}F_N(t)\,dt.
\end{align*}
補題\ref{lem:fejer}によって第2項は $N\to\infty$ で零へ向かう．
$x$ に関して一様な評価なので
$\limsup_N\norm{\sigma_Nf-f}_\infty\le\varepsilon$ を得る．
$\varepsilon$ は任意だから結論が従う．
\end{proof}

\begin{corollary}[三角多項式の稠密性]\label{cor:trig-dense}
連続な $2\pi$-周期関数全体の空間
$C(\mathbb T)$ において，
三角多項式全体は一様ノルムで稠密である．
\end{corollary}
\begin{proof}
$\sigma_Nf$ は次数高々 $N$ の三角多項式であり，
定理\ref{thm:fejer}から $\sigma_Nf\to f$ は一様収束する．
\end{proof}

\section{$L^2$ 完全性とParsevalの等式}

\begin{lemma}\label{lem:continuous-l2}
連続な周期関数全体 $C(\mathbb T)$ は
$L^2(\mathbb T)$ に稠密である．
\end{lemma}
\begin{proof}
まず単純関数は $L^2$ に稠密である．
任意の可測集合は，Lebesgue測度の正則性によって
測度に関して有限個の区間の和で近似できる．
各区間の指示関数は，その端点の小近傍で
値0から1へ線形に移行する連続周期関数によって，
任意に小さい $L^2$ 誤差で近似できる．
従って単純関数の有限線形結合も連続周期関数で
$L^2$ 近似できる．
三角不等式によって任意の $L^2$ 関数を近似できる．
\end{proof}

\begin{theorem}[Fourier級数の $L^2$ 完全性]\label{thm:l2}
任意の $f\in L^2(\mathbb T)$ について
\[
 \norm{S_Nf-f}_2\longrightarrow0,
\]
ただし $\norm f_2^2=(2\pi)^{-1}\int_{-\pi}^\pi|f|^2$
とする．さらに
\[
 \norm f_2^2=\sum_{n\in\mathbb Z}|\widehat f(n)|^2
\]
が成り立つ（Parsevalの等式）．
\end{theorem}
\begin{proof}
補題\ref{lem:continuous-l2}と
系\ref{cor:trig-dense}から，
三角多項式は $L^2(\mathbb T)$ に稠密である．
$S_Nf$ は正規直交系
$\{e^{int}:-N\le n\le N\}$ の張る部分空間への
直交射影である．
従って任意の次数高々 $N$ の三角多項式 $P$ に対して
\[
 \norm{f-S_Nf}_2\le\norm{f-P}_2
\]
が成り立つ．
任意の $\varepsilon>0$ に対して
三角多項式 $P$ で $\norm{f-P}_2<\varepsilon$ となるものを取る．
$N$ が $P$ の次数以上なら上の不等式から
$\norm{f-S_Nf}_2<\varepsilon$ となるので
$S_Nf\to f$ が $L^2$ で成立する．

直交性によるPythagorasの恒等式から
\[
 \norm f_2^2=\norm{S_Nf}_2^2+\norm{f-S_Nf}_2^2,
 \qquad
 \norm{S_Nf}_2^2=\sum_{n=-N}^N|\widehat f(n)|^2.
\]
$N\to\infty$ としてParsevalの等式を得る．
\end{proof}

\begin{corollary}[Fourier係数の一意性]\label{cor:fourier-unique}
$f,g\in L^2(\mathbb T)$ が
$\widehat f(n)=\widehat g(n)$ を全ての $n\in\mathbb Z$
について満たすなら，$f=g$ がほとんど至る所成立する．
\end{corollary}
\begin{proof}
$h=f-g$ と置くと全てのFourier係数は零である．
Parsevalの等式から
$\norm h_2^2=\sum_n|\widehat h(n)|^2=0$ なので
$h=0$ がほとんど至る所成立する．
\end{proof}

\section{各点収束}

\begin{theorem}[Dirichlet型収束定理]\label{thm:dirichlet}
$f$ を $2\pi$-周期の区分的 $C^1$ 関数とする．
各点 $x$ の左右極限を $f(x-),f(x+)$ と書くと
\[
 S_Nf(x)\longrightarrow\frac{f(x-)+f(x+)}2.
\]
\end{theorem}
\begin{proof}
Dirichlet核の積分表示から
\[
 S_Nf(x)=\frac1{2\pi}\int_{-\pi}^\pi
 f(x-t)\frac{\sin((N+\frac12)t)}{\sin(t/2)}\,dt.
\]
$t>0$ では $c_-=f(x-)$，
$t<0$ では $c_+=f(x+)$ と置き，
$c(t)=c_-$（$t>0$），$c_+$（$t<0$）と定義する．
$D_N$ は偶関数で
$(2\pi)^{-1}\int_{-\pi}^\pi D_N(t)\,dt=1$ なので
\[
 \frac1{2\pi}\int_{-\pi}^\pi c(t)D_N(t)\,dt
 =\frac{c_-+c_+}2.
\]
従って
\[
 S_Nf(x)-\frac{c_-+c_+}2
 =\frac1{2\pi}\int_{-\pi}^\pi
 \frac{f(x-t)-c(t)}{\sin(t/2)}
 \sin((N+\tfrac12)t)\,dt.
\]
$f$ が $x$ の両側で $C^1$ 級であることから，
$t\to0^\pm$ で
$f(x-t)-c(t)=O(|t|)$ となる．
従って
$g(t)=(f(x-t)-c(t))/\sin(t/2)$ は
$t=0$ の近傍で有界であり，
その外側でも区分的に連続だから
$g\in L^1([-\pi,\pi])$ である．
Riemann--Lebesgueの補題を
$g(t)e^{it/2}$ に適用すれば，
上の振動積分は $N\to\infty$ で零へ収束する．
よって結論が従う．
\end{proof}

\begin{corollary}[十分滑らかな場合の一様収束]\label{cor:uniform}
$f$ が $C^2$ 級の $2\pi$-周期関数なら，
Fourier級数は絶対かつ一様収束して $f$ に等しい．
\end{corollary}
\begin{proof}
$n\ne0$ に対して部分積分を二回行い，
周期性から端点の境界項が打ち消し合うので
\[
 \widehat f(n)=\frac1{(in)^2}\widehat{f''}(n).
\]
従って
$|\widehat f(n)|\le\norm{f''}_{L^1}/(2\pi n^2)$ である．
$\sum_{n\ne0}n^{-2}<\infty$ だから
WeierstrassのM判定法によりFourier級数は
絶対かつ一様収束する．
その和を $g$ とすると $g$ は連続であり，
項別積分から $\widehat g(n)=\widehat f(n)$ となる．
系\ref{cor:fourier-unique}と連続性から
$f=g$ が全ての点で成立する．
\end{proof}

\section{一様収束しない例と発散}

\begin{theorem}[連続関数のFourier部分和の発散]\label{thm:divergence}
ある連続な $2\pi$-周期関数 $f$ が存在して，
$S_Nf(0)$ が $N\to\infty$ で有界でない．
\end{theorem}
\begin{proof}
Banach空間 $C(\mathbb T)$ 上の連続線形汎関数
$T_Nf=S_Nf(0)$ を考える．
Dirichlet核の積分表示から
\[
 T_Nf=\frac1{2\pi}\int_{-\pi}^\pi f(-t)D_N(t)\,dt
\]
なので
\[
 \norm{T_N}\le\frac1{2\pi}\int_{-\pi}^\pi|D_N(t)|\,dt.
\]
逆向きの不等式を示す．
$\varepsilon>0$ に対して
$h_\varepsilon(t)=D_N(t)/(|D_N(t)|+\varepsilon)$ と置けば，
$h_\varepsilon$ は連続な実数値周期関数で
$\norm{h_\varepsilon}_\infty\le1$ である．
$f_\varepsilon(x)=h_\varepsilon(-x)$ と定めると
\[
 T_Nf_\varepsilon
 =\frac1{2\pi}\int_{-\pi}^\pi
   \frac{|D_N(t)|^2}{|D_N(t)|+\varepsilon}\,dt
 \longrightarrow\frac1{2\pi}\int|D_N(t)|\,dt
\]
が有界収束定理から従う．
従って
$\norm{T_N}=(2\pi)^{-1}\int|D_N|$ である．

$0<t\le\pi$ では
$|\sin(t/2)|\le t/2$ なので，
$A=N+\frac12$ と置けば
\begin{align*}
 \int_0^\pi|D_N(t)|\,dt
 &\ge2\int_0^\pi\frac{|\sin(At)|}{t}\,dt\\
 &=2\int_0^{A\pi}\frac{|\sin u|}{u}\,du\\
 &\ge2\sum_{k=1}^{N-1}
     \frac1{(k+1)\pi}\int_{k\pi}^{(k+1)\pi}|\sin u|\,du\\
 &=\frac4\pi\sum_{k=1}^{N-1}\frac1{k+1}
 \longrightarrow\infty.
\end{align*}
従って $\sup_N\norm{T_N}=\infty$ である．
一様有界性原理の対偶から，
ある $f\in C(\mathbb T)$ に対して
$\sup_N|T_Nf|=\infty$ となる．
すなわち $S_Nf(0)$ は有界でない．
\end{proof}

\begin{remark}
Fej\'erの定理は連続周期関数に対する
Ces\`aro平均の一様収束を保証するが，
通常のFourier部分和の各点収束を
すべての連続関数に対して保証するものではない．
一様有界性原理を用いることで，
この二つの収束現象の違いを明確に示せる．
\end{remark}


\section{到達課題の完全解答}
\noindent\textbf{鋸歯状関数のFourier級数．}
\begin{proposition}
$(-\pi,\pi)$ 上で $f(x)=x$ と定め，
$2\pi$ 周期拡張する．Fourier級数は
\[
f(x)\sim2\sum_{n=1}^\infty
\frac{(-1)^{n+1}}n\sin(nx)
\]
である．
\end{proposition}
\begin{proof}
$f$ は奇関数なので余弦係数は零である．
正弦係数は
\[
b_n=\frac2\pi\int_0^\pi x\sin(nx)\,dx
=\frac2\pi\left[
-\frac{x\cos(nx)}n+\frac{\sin(nx)}{n^2}
\right]_0^\pi
=\frac{2(-1)^{n+1}}n.
\]
従って表示された級数を得る．
区分的滑らかな周期関数のFourier収束定理から
$-\pi<x<\pi$ では $x$ に収束し，
不連続点では左右極限の平均0に収束する．
\end{proof}
\noindent\textbf{到達確認．}
$x=\pi/2$ を代入して交代奇数級数を求めよ．
\par\noindent\textbf{解答．}
$\pi/2=2(1-1/3+1/5-\cdots)$ なので
$1-1/3+1/5-\cdots=\pi/4$ である．
\end{document}